\documentclass[11pt,letterpaper]{article}

\usepackage[margin=1in]{geometry}
\usepackage{hyperref}
\usepackage{enumitem}
\usepackage{microtype}
\usepackage{parskip}
\usepackage{xcolor}
\usepackage{mdframed}
\usepackage{titlesec}

\hypersetup{
  colorlinks=true,
  linkcolor=blue!60!black,
  urlcolor=blue!60!black,
}

\title{\textbf{Brain Gene Expression Explorer}\\
  \large User Feedback \& Expert Evaluation}

\author{Thiago Viegas\\
  \small New York University}

\date{\today}

\begin{document}

\maketitle

\begin{abstract}
This document collects structured feedback on the \emph{Brain Gene Expression Explorer}
gathered from three domain experts during and after development.
The reviewers span neuroinformatics research, scientific software engineering, and
applied data science, providing complementary perspectives on both the scientific
utility and the design quality of the tool.
All three reviewers interacted with a live instance of the application.
\end{abstract}

% ─────────────────────────────────────────────────────────────────────────────
\section{Reviewer 1 — Ph.D.\ Researcher, NYU Neuroinformatics Lab}

\subsection*{Background}

This reviewer is a doctoral student in the same lab that develops the GTEx--AHBA
integration pipeline.
They contributed directly to the upstream modelling work and were a primary guide
throughout the development of the visualization, providing domain context, data
interpretation, and iterative feedback across multiple sessions.

\subsection*{Feedback on Scientific Utility}

The reviewer noted that the four-panel side-by-side layout immediately made
previously invisible patterns legible.
In particular, the simultaneous display of V1 (raw GTEx measurements) and V2
(LORO reconstructions) at the same held-out parcels allowed the team to verify,
at a glance, whether the model's predictions were tracking ground-truth values
or diverging systematically.

\begin{quote}
\itshape
``Being able to switch genes and immediately see how the prediction holds up
at the held-out parcels---without exporting CSVs or writing plot scripts---
has genuinely sped up how we diagnose model failures.
The residual map was especially useful: seeing the signed error in anatomical
space made it obvious that cerebellar parcels were consistently under-predicted
in a way that summary statistics alone never made clear.''
\end{quote}

The reviewer also noted that the V3 WBF-minus-AHBA-Reference residual map
(active for DLAM and PLAM) helped the team understand where the spatial
regularisation terms in those models were making a difference relative to
the naive baseline, providing visual evidence that the more complex models
were producing smoother and more anatomically plausible bilateral patterns.

\subsection*{Suggestions for Future Work}

Several of the future-work items in the main report originate directly from
this reviewer's suggestions:

\begin{enumerate}[nosep]
  \item \textbf{Expand subject coverage.}
        The reviewer observed that five donors is sufficient for early debugging
        but that differences in model behaviour across subjects—especially
        outlier donors—would only become visible with a larger cohort.
        They suggested adding support for additional GTEx donors as pipeline
        outputs become available.

  \item \textbf{Gene groups and functional categories.}
        Rather than exploring genes individually, the reviewer suggested
        allowing users to load predefined gene sets, such as pathway-annotated
        genes or high-differential-expression clusters, so that biological
        hypotheses can be interrogated across a curated subset rather than one
        gene at a time.

  \item \textbf{Uncertainty visualisation.}
        The upstream pipeline produces bootstrap confidence estimates alongside
        point predictions.
        The reviewer suggested exposing these as a secondary visual channel,
        such as parcel-level stippling or opacity modulation, to help
        distinguish high-confidence predictions from uncertain extrapolations.

  \item \textbf{Multi-subject overlay.}
        Averaging predictions across donors and displaying the group mean
        alongside individual subjects would help distinguish donor-specific
        idiosyncrasies from systematic model behaviour.
\end{enumerate}

% ─────────────────────────────────────────────────────────────────────────────
\section{Reviewer 2 — Software Engineer, Urban Data Analytics}

\subsection*{Background}

This reviewer is a software engineer with experience building interactive
geospatial and data visualization dashboards for urban analytics applications.
Although not a domain expert in neuroimaging, they brought a strong perspective
on interaction design, data density, and usability patterns common across
spatial visualization tools.

\subsection*{Feedback on Interaction Design}

The reviewer evaluated the tool from the perspective of a technically
sophisticated user encountering it for the first time without domain
background.
Their primary observation was that the colour-only encoding of the 3D brain
panels, while visually engaging, made it difficult to recover the identity and
precise value of a parcel without additional context.

\begin{quote}
\itshape
``When you have a spatial view coloured by a continuous variable, the natural
next question is always: which region is that, and what is its exact value?
Without hover or click feedback, the brain looks beautiful but you are
essentially reading a map with no legend for the individual features.''
\end{quote}

Based on this feedback, two features were implemented:

\begin{enumerate}[nosep]
  \item \textbf{Parcel-level click tooltips.}
        Clicking any parcel in any of the four panels now displays a floating
        card showing the parcel name and expression value.
        In residual mode, the tooltip expands to show Truth, Predicted, and
        $\Delta$ for V2, and AHBA Reference, WBF, and $\Delta$ for V3,
        giving exact numerical context for whatever is visible in colour.

  \item \textbf{V5 Model Analytics panel.}
        The reviewer suggested adding a supplementary statistical view that
        would allow users to identify where each parcel sits in the overall
        distribution of predictions---not just its colour on the surface.
        This motivated the V5 panel: a two-row analytics overlay with a
        colour-coded held-out parcel scatter (row 1, current gene) and a
        global per-gene mean-expression scatter alongside a ranked
        Pearson-$r$ bar chart (row 2, all genes).
        The reviewer noted that position and length are more precise
        quantitative channels than colour, and that the combination allows
        users to contextualise any individual parcel or gene within the
        full dataset without leaving the application.
\end{enumerate}

The reviewer also praised the opacity sliders for cortex, subcortex, and
cerebellum as an effective solution to the occlusion problem that is typical
in 3D spatial views, drawing a parallel to layer-toggle controls in GIS
mapping tools.

% ─────────────────────────────────────────────────────────────────────────────
\section{Reviewer 3 — Professor of Neuroinformatics}

\subsection*{Background}

This reviewer is a faculty member in neuroinformatics whose research involves
computational methods for integrating multimodal brain data, including
transcriptomic and imaging datasets of the type used in this project.
They evaluated the tool as an external expert in both the domain and the
broader landscape of neuroimaging visualization.

\subsection*{Endorsement and Key Observations}

The reviewer endorsed the tool as a meaningful contribution to the model
development workflow, emphasising in particular its value for understanding
how the models extrapolate sparse measurements to the rest of the brain.

\begin{quote}
\itshape
``What strikes me about this tool is that it makes the extrapolation step
legible. The models take a handful of measured parcels and predict expression
everywhere else---that is a significant inferential leap, and V3 makes the
result of that leap visible in anatomical space. Being able to compare the
dense whole-brain fit against the sparse ground truth in the same frame of
reference, and then toggle to a residual view, is exactly the kind of check
that should happen before these predictions are used for downstream analysis.''
\end{quote}

The reviewer highlighted two specific aspects as particularly well-executed:

\begin{enumerate}[nosep]
  \item \textbf{Anatomical faithfulness.}
        Mapping predictions directly onto the parcellation atlas, with correct
        subcortical mesh geometry and hemisphere handling, ensures that the
        spatial interpretation of the visualisation is biologically grounded
        rather than schematic.
        The reviewer noted that many model-debugging dashboards use simplified
        region representations that can obscure anatomically specific patterns.

  \item \textbf{Model-variant comparison.}
        The ability to switch between naive, DLAM, and PLAM within the same
        four-panel layout, with the naive model fixed as the V4 reference,
        was identified as a strong design choice that supports rigorous
        comparative evaluation without requiring users to maintain multiple
        browser tabs or coordinate separate export steps.
\end{enumerate}

The reviewer also noted that the tool would be well-suited for use in a
graduate seminar context, where trainees need to develop intuition for how
spatial transcriptomic models behave before interpreting their outputs
quantitatively.

% ─────────────────────────────────────────────────────────────────────────────
\section{Summary}

\begin{center}
\begin{tabular}{p{3.2cm} p{3.5cm} p{6.5cm}}
\hline
\textbf{Reviewer} & \textbf{Background} & \textbf{Primary contribution} \\
\hline
Ph.D.\ researcher & Neuroinformatics, pipeline co-developer &
  Validated scientific utility; sourced future-work suggestions including
  subject expansion, gene groups, and uncertainty visualisation \\[4pt]
Software engineer & Urban data / geospatial dashboards &
  Identified colour-only encoding limitation; motivated click tooltips
  and V5 analytics panel \\[4pt]
Neuroinformatics professor & Faculty, multimodal brain data &
  Endorsed extrapolation legibility and anatomical faithfulness; suggested
  pedagogical use case \\
\hline
\end{tabular}
\end{center}

\end{document}
